%%
%% Capítulo 5: Conclusões
%%

\mychapter{Conclusão}
\label{Cap:conclusao}

Este trabalho propôs a aplicação de uma arquitetura ConvLSTM para a predição de crimes violentos contra o patrimônio na cidade de Maceió, utilizando dados georreferenciados de ocorrências entre 2012 e 2022. O pipeline desenvolvido abrangeu desde o pré-processamento dos dados brutos (limpeza, filtragem espacial por polígono municipal, agrupamento de tipos de crime) até a construção de representações espaço-temporais em grids, o treinamento do modelo e a avaliação com métricas tradicionais e com tolerância espacial.

Os resultados preliminares obtidos para Maceió, com semente aleatória 42, indicam que o modelo ConvLSTM é capaz de identificar regiões com maior probabilidade de ocorrência criminal. No threshold de 0,50, o modelo apresentou recall de 72,8\% e precision de 5,9\%, resultando em um F1 de 10,5\%. Esses valores refletem o desafio inerente da esparsidade dos dados: como apenas cerca de 2\% das células apresentam crime em um dado dia, a taxa de falsos positivos é elevada. 

Ao se considerar a avaliação com tolerância espacial (vizinhança de Chebyshev $\leq 1$), os resultados apresentaram melhora expressiva: recall espacial de 86,4\%, precision espacial de 33,1\% e F1 espacial de 45,4\% no threshold de 0,50. O melhor F1 espacial observado foi de 47,8\% no threshold de 0,65, com recall espacial de 77,7\% e precision espacial de 37,5\%. Esses resultados indicam que, embora o modelo não acerte exatamente a célula, uma parcela significativa das previsões incide em vizinhanças imediatas de crimes reais, o que pode ter utilidade prática para o direcionamento de ações preventivas.

Ao longo do treinamento de 200 épocas, observou-se uma redução consistente tanto na perda de treinamento (de 1,358 para 0,743) quanto na perda de validação (de 1,102 para 0,697), sugerindo que o modelo não apresentou sobreajuste significativo apesar do alto dropout aplicado ($p = 0{,}8$).

É importante destacar algumas limitações do presente estudo. A precisão tradicional permaneceu abaixo de 9\% em todos os thresholds, o que significa que a maioria dos alertas gerados pelo modelo não correspondem a ocorrências reais na célula exata. Esse padrão é consistente com os achados da literatura \cite{albors_zumel_deep_2025}, que reporta limitações análogas em cidades norte-americanas. Além disso, o pipeline atual utiliza apenas o canal de crime (modo ``C''), sem incorporar dados sociodemográficos ou de mobilidade, que poderiam contribuir para a melhora do desempenho conforme sugerido por \citeasnoun{albors_zumel_deep_2025}.

O pipeline para Arapiraca encontra-se em fase de pré-processamento: os dados foram filtrados espacialmente (19.141 registros dentro do polígono municipal), mas as etapas de construção de grids, preparação das sequências e treinamento do modelo ainda necessitam ser concluídas.


\section{Trabalhos Futuros}
\label{Sec:TrabalhosFuturos}

Com base nas limitações identificadas e nas possibilidades abertas pelo pipeline desenvolvido, os seguintes trabalhos futuros podem ser realizados:

\begin{itemize}
  \item Completar o pipeline para a cidade de Arapiraca, incluindo a construção de grids, treinamento e avaliação do modelo ConvLSTM, permitindo a comparação entre as duas cidades;
  \item Incorporar canais adicionais de entrada (dados sociodemográficos e de mobilidade) para avaliar se informações contextuais melhoram o desempenho preditivo, conforme o modo ``CMS'' proposto na literatura de referência;
  \item Experimentar diferentes granularidades temporais (12h, 8h) e tamanhos de janela de lookback (14 dias), comparando o impacto na qualidade da predição;
  \item Implementar e comparar modelos baseline (persistência do último mapa, média móvel, regressão logística, Random Forest) para quantificar o ganho efetivo da ConvLSTM;
  \item Avaliar o impacto de diferentes tamanhos de grid e de sub-grid na capacidade preditiva do modelo;
  \item Investigar técnicas de aumento de dados e estratégias alternativas de balanceamento (focal loss, oversampling de amostras positivas) para mitigar o impacto da esparsidade;
  \item Explorar a generalização do modelo por meio de transferência de aprendizado entre cidades (treinamento em Maceió e avaliação em Arapiraca, e vice-versa).
\end{itemize}

